%%%%%%%%%%%%%%%%%%%%%%%%%%%%%%%%%%%%%%%%%%%%%%%%%%%%%%%%%%%%%%%%%%%%%%%%%%%%%%%
% Riccardo Gonzo - Academic CV
% Default and master versions, revised July 2026
%%%%%%%%%%%%%%%%%%%%%%%%%%%%%%%%%%%%%%%%%%%%%%%%%%%%%%%%%%%%%%%%%%%%%%%%%%%%%%%

\PassOptionsToPackage{colorlinks=true,pdfpagelabels=false}{hyperref}
\documentclass[11pt,a4paper,sans]{moderncv}
\moderncvstyle[details,left,nosymbols]{classic}
\moderncvcolor{blue}

\usepackage[T1]{fontenc}
\usepackage[utf8]{inputenc}
\usepackage{lmodern}
\usepackage[margin=0.67in,top=0.56in,bottom=0.66in,footskip=0.30in]{geometry}
\usepackage{microtype}
\usepackage{ragged2e}
\usepackage{fancyhdr}
\usepackage{lastpage}
\usepackage{hyperref}
\usepackage{xparse}
\usepackage[
  backend=biber,
  style=numeric,
  sorting=none,
  defernumbers=true,
  maxbibnames=99,
  giveninits=true,
  doi=false,
  eprint=false,
  url=false,
  isbn=false
]{biblatex}
\addbibresource{Riccardo_Gonzo_Publications_Final_v4.bib}

% Restrained palette: dark ink for hierarchy, warm ochre for emphasis,
% muted blue for links, and grey for dates/secondary information.
\definecolor{CVInk}{HTML}{263945}
\definecolor{CVOchre}{HTML}{A96F16}
\definecolor{CVOchreLight}{HTML}{D7C29A}
\definecolor{CVLink}{HTML}{2C6077}
\definecolor{CVText}{HTML}{262626}
\definecolor{CVGrey}{HTML}{666C70}

\colorlet{color0}{CVText}
\colorlet{color1}{CVInk}
\colorlet{color2}{CVGrey}
\colorlet{firstnamecolor}{CVInk}
\colorlet{lastnamecolor}{CVInk}
\colorlet{namecolor}{CVInk}
\colorlet{titlecolor}{CVOchre}
\colorlet{addresscolor}{CVLink}
\colorlet{sectioncolor}{CVInk}
\colorlet{subsectioncolor}{CVInk}
\colorlet{bodyrulecolor}{CVOchreLight}
\colorlet{hintstylecolor}{CVGrey}
\colorlet{headrulecolor}{CVOchreLight}
\colorlet{quotecolor}{CVOchre}

\hypersetup{
  urlcolor=CVLink,
  linkcolor=CVLink,
  citecolor=CVLink,
  pdfauthor={Riccardo Gonzo},
  pdfsubject={Academic curriculum vitae},
  pdfborder={0 0 0},
  bookmarksopen=true
}
\urlstyle{same}

\renewcommand*{\namefont}{\fontsize{31}{33}\mdseries\upshape}
\renewcommand*{\titlefont}{\fontsize{12}{14}\bfseries\upshape}
\renewcommand*{\addressfont}{\fontsize{9.2}{11.2}\selectfont\mdseries\upshape}
\renewcommand*{\sectionfont}{\fontsize{14.2}{16.2}\selectfont\bfseries\upshape}
\renewcommand*{\subsectionfont}{\fontsize{11.6}{13.2}\selectfont\bfseries\upshape}
\renewcommand*{\hintfont}{\fontsize{9.1}{10.5}\selectfont\bfseries\upshape}

% The dates are large enough to read, but the gutter remains compact.
\setlength{\hintscolumnwidth}{1.82cm}
\setlength{\separatorcolumnwidth}{0.10cm}
\setlength{\parindent}{0pt}
\setlength{\parskip}{0pt}
\emergencystretch=2em
\raggedbottom

\makeatletter
\RenewDocumentCommand{\section}{sm}{%
  \par\addvspace{1.48ex}%
  \phantomsection%
  \addcontentsline{toc}{section}{#2}%
  \noindent\sectionstyle{#2}\par\nobreak
  \vspace{-0.25ex}%
  \noindent{\color{CVOchreLight}\rule{\textwidth}{0.62pt}}\par\nobreak
  \vspace{0.43ex}\@afterheading}
\RenewDocumentCommand{\subsection}{sm}{%
  \par\addvspace{1.02ex}%
  \phantomsection%
  \addcontentsline{toc}{subsection}{#2}%
  \noindent\textcolor{CVOchre}{\rule{1.8pt}{1.02em}}\hspace{0.46em}%
  \subsectionstyle{#2}\par\nobreak
  \vspace{0.22ex}\@afterheading}
\makeatother

\newcommand{\talkitem}[4]{\cvitem[0.11em]{#1}{\textbf{#2}: \emph{#3}#4}}
\newcommand{\linktag}[2]{\href{#1}{\textcolor{CVLink}{[#2]}}}
\newcommand{\referenceitem}[3]{%
  \par\noindent\textbf{#1}, #2. \href{mailto:#3}{#3}\par\vspace{0.14em}%
}

\DeclareBibliographyCategory{preprints}
\addtocategory{preprints}{Gonzo:2026yha,Alessio:2026bdi}

\DeclareFieldFormat{title}{#1}
\DeclareFieldFormat{journaltitle}{\mkbibemph{#1}}
\DeclareFieldFormat{volume}{\mkbibbold{#1}}
\DeclareFieldFormat{pages}{#1}
\DeclareFieldFormat{doi}{\href{https://doi.org/#1}{\textcolor{CVLink}{[DOI]}}}
\DeclareFieldFormat{eprint:arxiv}{\href{https://arxiv.org/abs/#1}{\textcolor{CVLink}{[arXiv:#1]}}}
\AtEveryBibitem{%
  \clearfield{number}%
  \clearfield{month}%
  \clearfield{reportnumber}%
  \clearfield{primaryclass}%
}
\DeclareBibliographyDriver{article}{%
  \usebibmacro{bibindex}%
  \usebibmacro{begentry}%
  \printnames{author}%
  \setunit{\addperiod\space}%
  \printfield{title}%
  \setunit{\addperiod\space}%
  \iffieldundef{journaltitle}
    {\printtext{Preprint}\setunit{\addcomma\space}\printfield{year}}
    {\printfield{journaltitle}%
     \setunit{\addspace}%
     \printfield{volume}%
     \iffieldundef{pages}{}{\setunit{\addcomma\space}\printfield{pages}}%
     \setunit{\addspace}%
     \printtext{\mkbibparens{\printfield{year}}}}%
  \iffieldundef{note}{}{\setunit{\addperiod\space}\printfield{note}}%
  \iffieldundef{doi}{}{\setunit{\addspace}\printfield{doi}}%
  \iffieldundef{eprint}{}{\setunit{\addspace}\printfield[eprint:arxiv]{eprint}}%
  \usebibmacro{finentry}}

\defbibenvironment{bibliography}
  {\list
     {\printtext[labelnumberwidth]{\printfield{labelnumber}}}
     {\setlength{\topsep}{0pt}%
      \setlength{\partopsep}{0pt}%
      \setlength{\itemsep}{1.15pt}%
      \setlength{\parsep}{0pt}%
      \setlength{\labelwidth}{0.48cm}%
      \setlength{\labelsep}{0.18cm}%
      \setlength{\leftmargin}{0.66cm}%
      \setlength{\itemindent}{0pt}%
      \renewcommand*{\makelabel}[1]{\hss##1}}}
  {\endlist}
  {\item}

\fancypagestyle{cvfooter}{%
  \fancyhf{}%
  \renewcommand{\headrulewidth}{0pt}%
  \renewcommand{\footrulewidth}{0pt}%
  \fancyfoot[L]{\scriptsize\textcolor{CVGrey}{Riccardo Gonzo -- Curriculum Vitae -- updated July 2026}}%
  \fancyfoot[R]{\scriptsize\textcolor{CVGrey}{\thepage/\pageref*{LastPage}}}}
\AtBeginDocument{\pagestyle{cvfooter}}

\firstname{Riccardo}
\familyname{Gonzo}
\title{Royal Society University Research Fellow}

\newcommand{\cvheader}{%
  \noindent
  \begin{minipage}[t]{0.59\textwidth}
    \vspace{0pt}%
    {\fontsize{31}{33}\selectfont\color{CVInk}Riccardo Gonzo}\par
    \vspace{0.22em}%
    {\fontsize{12}{14}\selectfont\bfseries\color{CVOchre}Royal Society University Research Fellow}
  \end{minipage}\hfill
  \begin{minipage}[t]{0.38\textwidth}
    \vspace{0pt}%
    \raggedleft\fontsize{9.2}{11.2}\selectfont
    \textbf{Email:} \href{mailto:r.gonzo@qmul.ac.uk}{r.gonzo@qmul.ac.uk}\\
    \textbf{Website:} \href{https://riccardogonzo.com}{riccardogonzo.com}\\
    \textbf{ORCID:} \href{https://orcid.org/0000-0001-7285-6295}{0000-0001-7285-6295}\\
    \textbf{INSPIRE--HEP:} \href{https://inspirehep.net/authors/1812058?ui-citation-summary=true}{author profile}
  \end{minipage}%
  \par\vspace{0.78em}%
}

\begin{document}
\cvheader
\thispagestyle{cvfooter}

\section{Research profile}
I am a theoretical physicist with over eight years of research experience spanning high-energy physics and general relativity. My work focuses on the analytical description of gravitational binary systems and their waveforms, combining particle-physics tools such as scattering amplitudes with post-Newtonian and gravitational self-force methods. During my PhD, I developed a direct link between on-shell scattering amplitudes and gravitational waveforms. Since then, I have worked on strong-field scatter-to-bound maps for Kerr orbits and waveforms, and on an effective-field-theory formulation of the self-force. I have developed this work in a highly collaborative international environment, working with more than thirty researchers across Europe, North America and Asia. In 2026, I was awarded the prestigious \href{https://royalsociety.org/grants/university-research/}{Royal Society University Research Fellowship}, a major eight-year fellowship with a total programme value of approximately \pounds{}1.8 million, at the University of Southampton for the project \emph{High-precision effective field theory for extreme-mass-ratio inspirals}.

My long-term vision is to create a unified framework that connects weak-field Post-Minkowskian and post-Newtonian theory with strong-field methods from the self-force approach, leading to improved waveform models for current (LIGO--Virgo--KAGRA) and future (LISA, ET and CE) gravitational-wave detectors. At Southampton I will build an independent group with two PhD students and a postdoctoral researcher, working on particle-physics methods for high-precision gravitational-wave modelling, with a particular focus on LISA.

\section{Academic appointments}
\cventry{2026--}{Royal Society University Research Fellow}{University of Southampton}{Southampton}{}{}
\cventry{2025--2026}{Postdoctoral Research Associate}{Queen Mary University of London}{London}{}{}
\cventry{2022--2025}{Postdoctoral Research Associate}{University of Edinburgh}{Edinburgh}{}{}

\section{Fellowships and research funding}
\cvitem[0.24em]{2026}{\textbf{Royal Society University Research Fellowship}. Major eight-year fellowship; total programme value approximately \pounds{}1.8 million. Host: University of Southampton. Project: \emph{High-precision effective field theory for extreme-mass-ratio inspirals}. \linktag{https://royalsociety.org/grants/university-research/}{scheme}}
\cvitem[0.24em]{2019}{\textbf{Marie Sk\l odowska-Curie ITN fellowship, SAGEX}. European Union Horizon 2020 fellowship, providing approximately \texteuro100k of personal fellowship support and advanced research and professional training over three years. \linktag{https://sagex.ph.qmul.ac.uk/}{SAGEX}}
\cvitem[0.24em]{2023--2026}{\textbf{\pounds{}139.7k secured in competitive funding for international programmes and workshops}: GGI six-week programme (\pounds{}83.7k); ICTP-SAIFR two-week programme (\pounds{}25k); Nordita two-week programme (\pounds{}20.6k); and Higgs Centre workshop (\pounds{}10.4k). Co-applicant and organiser for all four; main coordinator for the GGI, Nordita and Higgs Centre activities. Details are given below.}

\section{Education}
\cventry{2018--2022}{PhD in Theoretical Particle Physics}{Trinity College Dublin}{Dublin}{}{Supervisor: Prof. Ruth Britto. Thesis: \emph{Coherent states and classical radiative observables in the S-matrix formalism} \linktag{http://www.tara.tcd.ie/handle/2262/98491}{thesis}. During my PhD, I won a prestigious and competitive Marie-Curie International Training Network fellowship, \emph{Scattering Amplitudes: from Geometry to Experiments} (SAGEX). The fellowship gave me research training through dedicated workshops and schools and a three-month internship at Wolfram Research, as well as training in scientific writing, presentation skills, STEM outreach, grant writing, time and project management, career development and good scientific practice. A detailed record is available on the \href{https://sagex.org/}{SAGEX website}.}
\cventry{2015--2017}{MSc in Theoretical Physics}{University of Padova}{Padova}{\textit{110/110 cum laude}}{Erasmus+ exchange at the University of G\"ottingen for the MSc thesis \emph{The infinite-spin representations of the Poincar\'e group}. Additional coursework: five examinations beyond the degree requirements (34 ECTS).}
\cventry{2012--2015}{BSc in Physics}{University of Padova}{Padova}{\textit{110/110 cum laude}}{Additional coursework: one examination beyond the degree requirements (6 ECTS).}

\clearpage
\section{Publications and research impact}
My full publication list is available through my \href{https://inspirehep.net/authors/1812058?ui-citation-summary=true}{INSPIRE--HEP profile} and \href{https://orcid.org/0000-0001-7285-6295}{ORCID record}. I have 22 papers, including 20 published in renowned journals such as \emph{Journal of High Energy Physics}, \emph{Physical Review D} and \emph{Physical Review Letters}, with over 1000 citations and an $h$-index of 16 according to INSPIRE--HEP (July 2026).

\begin{refsection}
\nocite{Gonzo:2026yha,Alessio:2026bdi}
\subsection{Preprints}
{\footnotesize\printbibliography[heading=none]}
\end{refsection}

\begin{refsection}
\nocite{Barack:2026izc,Alessio:2025isu,Akpinar:2025tct,Alessio:2025flu,Akpinar:2025huz,Gonzo:2024xjk,Gonzo:2024zxo,Adamo:2024oxy,DeAngelis:2023lvf,Cristofoli:2021jas,Gonzo:2023cnv,Gonzo:2023goe,Adamo:2022ooq,Gonzo:2022tjm,Britto:2021pud,Cristofoli:2021vyo,Gonzo:2021drq,Gonzo:2020xza,Gonzo:2019fai,Banerjee:2020kaa}
\subsection{Peer-reviewed publications}
{\footnotesize\printbibliography[heading=none]}
\end{refsection}

\section{International programmes and workshop series}
In parallel with my research, I started a series of meetings to connect the general-relativity and particle-physics approaches to binary dynamics, which had often developed separately. The first workshop in Edinburgh focused on gravitational self-force and scattering amplitudes; the later meetings broadened the discussion to post-Newtonian, post-Minkowskian, effective-one-body, numerical-relativity and data-analysis methods. The aim has been practical: to compare methods, identify calculations that can be cross-checked, and start collaborations on waveform modelling, especially for LISA. I secured competitive funding for the series and was the main coordinator for Edinburgh, Nordita and the forthcoming GGI programme.

\cvitem[0.22em]{2027}{\textbf{Main coordinator and co-organiser}, forthcoming GGI programme \emph{Synergies for Black-Hole Binaries: from Perturbative Techniques to Strong-Field Gravity}, Florence. Six-week programme including a conference week.}
\cvitem[0.22em]{2027}{\textbf{Co-organiser}, forthcoming ICTP-SAIFR / Principia programme \emph{All Roads Lead to Waveforms: GSF, PM, PN and All That}, S\~ao Paulo. Two-week programme including a PhD school.}
\cvitem[0.22em]{2026}{\textbf{Main coordinator and co-organiser}, \emph{Amplitudes, Strong-Field Gravity, and Resummation}, Nordita, Stockholm. Two-week programme including a one-week PhD school and a one-week workshop, with about 120 participants. \linktag{https://indico.fysik.su.se/event/9143/overview}{programme}}
\cvitem[0.22em]{2025}{\textbf{Co-organiser}, \emph{2nd Annual Workshop on Self-Force and Amplitudes}, University of Southampton. One-week workshop with about 70 participants. \linktag{https://indico.cern.ch/event/1485758/}{workshop}}
\cvitem[0.22em]{2024}{\textbf{Main coordinator and co-organiser}, \emph{Gravitational Self-Force and Scattering Amplitudes}, Higgs Centre, Edinburgh. One-week workshop with about 40 participants. \linktag{https://higgs.ph.ed.ac.uk/workshops/gravitational-self-force-and-scattering-amplitudes/}{workshop}}

\subsection{Other scientific organisation}
\cvitem[0.22em]{2025--2026}{\textbf{Local organiser}, QMUL Centre for Theoretical Physics journal club; coordinated regular biweekly meetings.}
\cvitem[0.22em]{2022--2023}{\textbf{Co-organiser}, GRAMPA online seminar series on scattering amplitudes and their applications to gravitational waves.}
\cvitem[0.22em]{2021}{\textbf{Co-organiser}, SAGEX Amplitudes School, Niels Bohr Institute, held in connection with the Amplitudes 2021 conference. I helped coordinate the lecturers and the online school programme. \linktag{https://indico.nbi.ku.dk/event/1530/}{school}}

\section{Competitive programme and workshop funding}
\cvitem[0.18em]{2026}{\textbf{\pounds{}83.7k}, GGI programme \emph{Synergies for Black-Hole Binaries: from Perturbative Techniques to Strong-Field Gravity}; co-applicant and main coordinator.}
\cvitem[0.18em]{2026}{\textbf{\pounds{}25k}, ICTP-SAIFR / Principia programme \emph{All Roads Lead to Waveforms: GSF, PM, PN and All That}; co-applicant and organiser.}
\cvitem[0.18em]{2025}{\textbf{\pounds{}20.6k}, Nordita programme \emph{Amplitudes, Strong-Field Gravity, and Resummation}; co-applicant and main coordinator.}
\cvitem[0.18em]{2023}{\textbf{\pounds{}10.4k}, Higgs Centre workshop \emph{Gravitational Self-Force and Scattering Amplitudes}; co-applicant and main coordinator.}

\section{Supervision and mentoring}
\cventry{2024--}{PhD co-supervisor: Emanuele Rosi}{Sapienza University of Rome}{}{}{Co-supervision with Prof. Vittorio Del Duca. I proposed the initial project on gravitational waves and scattering amplitudes for compact binaries. We meet weekly, and the collaboration has already led to two joint papers: \linktag{https://arxiv.org/abs/2511.11457}{arXiv:2511.11457} and \linktag{https://arxiv.org/abs/2601.21687}{arXiv:2601.21687}.}
\cventry{2024}{MSc co-supervisor: Andrea de Simone}{Sapienza University of Rome}{}{}{Six-month thesis on gravitational scattering amplitudes and the two-body problem; co-supervision with Prof. Vittorio Del Duca.}
\cventry{2024}{MSc co-supervisor: Damiano Barcaro}{Sapienza University of Rome}{}{}{Six-month thesis: \emph{Scattering waveforms for Kerr black holes from the soft expansion} \linktag{https://arxiv.org/abs/2411.18632}{paper}; subsequently began a PhD at Johannes Gutenberg University Mainz.}
\cventry{2022}{MSc supervisor: Leixi Wang}{University of Edinburgh}{}{}{Four-month thesis, \emph{Classical spacetimes from amplitudes}; I proposed the project and supervised it to completion.}

\section{Invited talks, discussion leadership and seminars}
\subsection{Invited talks at conferences and workshops}
{\small
\talkitem{May 2026}{CERN, Gravitational Waves from Amplitudes}{Gravitational scattering in the ultra-relativistic limit}{, Geneva.}
\talkitem{Mar. 2026}{1st APCTP--London Triangle Global Collaboration Workshop}{What particle theory teaches us about binary black holes}{, King's College London.}
\talkitem{Mar. 2026}{Black Hole Perturbations and Holography, IGAP}{Self-force EFT for black-hole binaries: bridging weak- and strong-field expansions via critical orbits}{, Trieste. \linktag{https://www.youtube.com/watch?v=oDev489dbvo}{video}}
\talkitem{Mar. 2026}{Max Planck Group Leader Symposium}{Particle theory for gravitational waveform modelling}{, Harnack Haus, Berlin.}
\talkitem{Nov. 2025}{FPUK Meeting 2025}{Unexpected symmetries and on-shell integrability of Kerr black-hole scattering}{, King's College London. \linktag{https://www.youtube.com/watch?v=MNupVDsL_4c}{video}}
\talkitem{June 2025}{Amplitudes 2025}{Dirac brackets for classical black-hole scattering: from amplitudes to observables}{, Seoul. \linktag{https://www.youtube.com/watch?v=r0f0wpcV9S4}{video}}
\talkitem{May 2025}{Investigating Compact Objects and Binary Mergers with Gravitational Waves}{Worldline approach for Kerr black holes: PM vs GSF}{, Sapienza University of Rome.}
\talkitem{Feb. 2025}{First Quantisation for Physics in Strong Fields}{Worldline approach for spinning bound binaries}{, Higgs Centre, Edinburgh.}
\talkitem{Feb. 2025}{String Theory as a Bridge between Gauge Theory and Quantum Gravity}{S-matrix tools for bound waveform modelling}{, Sapienza University of Rome.}
\talkitem{Jan. 2025}{KITP, What Is Particle Theory?}{Classical bound observables from the S-matrix formalism}{, Santa Barbara. \linktag{https://online.kitp.ucsb.edu/online/particles25/gonzo/rm/jwvideo.HTML}{video}}
\talkitem{Oct. 2024}{MIAPbP}{On-shell maps between scattering and bound observables}{, Munich.}
\talkitem{Sept. 2024}{AEI, Fundamental Physics Meets Waveforms with LISA}{Black holes as point particles: from amplitudes to bound waveforms}{, Potsdam.}
\talkitem{July 2024}{ECCF 2024}{Light-ray operators in gauge and gravity theory: new infrared-finite local observables}{, Mainz.}
\talkitem{Feb. 2024}{JENAS Initiative: Gravitational Wave Probes of Fundamental Physics}{Black holes as point particles: from amplitudes to self-force}{, Rome.}
\talkitem{Jan. 2024}{GRAMPA / ICMS}{Classical gravitational bound states with amplitudes}{, Edinburgh. \linktag{https://www.youtube.com/watch?v=ehdMRQqH4IQ&list=PLUbgZHsSoMEV_4j8ratb67TXYiHxURZ_c&index=7}{video}}
\talkitem{Aug. 2023}{GWs Meet Amplitudes}{From classical scattering amplitudes to bound-state observables}{, S\~ao Paulo. \linktag{https://www.youtube.com/watch?v=_z1hLS65mqA}{video}}
\talkitem{Dec. 2022}{QCD Meets Gravity}{Bethe--Salpeter equation for classical gravitational bound states}{, Z\"urich. \linktag{https://indico.phys.ethz.ch/event/22/contributions/197/attachments/150/266/Gonzo_Video.mp4}{video}}
\talkitem{Sept. 2022}{Celestial Amplitudes and Flat-Space Holography}{Celestial holography on non-trivial backgrounds}{, Corfu. \linktag{https://www.youtube.com/watch?v=QltnFM1-HCo}{video}}
\talkitem{Aug. 2022}{Physics in Intense Fields, PIF22}{High-energy limit of quantum and classical wave-scattering observables}{}
\talkitem{Dec. 2021}{QCD Meets Gravity}{An infinity of amplitude relations in classical physics}{, UCLA. \linktag{https://www.youtube.com/watch?v=KnxY-lWn9u0}{video}}
\talkitem{May 2021}{GGI workshop}{Waveforms from the KMOC formalism and coherent states}{, Florence. \linktag{https://www.youtube.com/watch?v=ZYaCQPMhyEE&t=268s}{video}}
}

\subsection{Contributed talks}
{\small
\talkitem{Sept. 2026}{Avenues of Quantum Field Theory in Curved Spacetime 2026}{Weak-to-strong field resummation in black-hole binaries}{, Lisbon. \linktag{https://inspirehep.net/conferences/3153477}{event}}
\talkitem{July 2026}{Capra Conference 2026}{Self-force EFT for binary dynamics: a worldline approach to bound motion}{, Brussels.}
\talkitem{July 2025}{Capra Conference 2025}{Worldline approach to Kerr black holes: PM vs GSF}{, Southampton.}
\talkitem{Apr. 2025}{BritGrav 2025}{Worldline approach to Kerr black holes: PM vs GSF}{, University of Birmingham.}
\talkitem{Sept. 2023}{New Frontiers in Theoretical Physics}{Classical gravitational bound states with amplitudes}{, Cortona.}
\talkitem{Dec. 2020}{XVI Avogadro Meeting}{Light-ray operators and detector algebra}{}
}

\subsection{Discussion panels and session leadership}
{\small
\talkitem{July 2025}{Capra Conference 2025}{Invited discussion panellist, ``Synergies and hybrid models''}{, Southampton.}
\talkitem{June 2024}{S-matrix Bootstrap 2024}{Invited talk for the discussion session ``Gravitational waves and scattering amplitudes''}{, Reykjav\'ik.}
\talkitem{Mar. 2024}{Gravitational Self-Force and Scattering Amplitudes}{Leader of the discussion session ``Self-force and amplitudes''}{, Edinburgh.}
\talkitem{Sept. 2022}{Celestial Amplitudes and Flat-Space Holography}{Invited lecture for the discussion session ``Connections between gravity, classical observables and scattering amplitudes''}{, Corfu.}
}

\subsection{Invited seminars}
{\small
\talkitem{June 2026}{CERN}{High-energy scattering in gravity: Regge theory, shock waves and beyond}{, Geneva.}
\talkitem{Mar. 2026}{GRAMPA seminar series}{Self-force EFT: weak-to-strong-field resummation in black-hole binaries}{, online.}
\talkitem{Dec. 2025}{IHES}{Regge theory for gravity amplitudes}{, Paris.}
\talkitem{Nov. 2025}{Amplitudes Lounge Seminar}{High-energy gravitational scattering: Regge theory and shock-wave formalism}{, online.}
\talkitem{June 2025}{AEI}{Dirac brackets for classical black-hole scattering}{, Berlin.}
\talkitem{Feb. 2025}{UCLA}{Insights into the scattering-to-bound map from the self-force approach}{}
\talkitem{Nov. 2024}{NBI}{The on-shell space of scattering and bound classical observables}{, Copenhagen.}
\talkitem{Sept. 2024}{Humboldt University DFG retreat}{Shedding light on bound states with the S-matrix}{, Berlin.}
\talkitem{Sept. 2024}{Worldline Seminars}{Classical scattering and bound observables from the worldline approach to the two-body problem}{, online.}
\talkitem{May 2024}{University of Z\"urich}{From amplitudes to gravitational bound observables}{}
\talkitem{May 2024}{Scuola Normale Superiore}{Classical bound observables from amplitudes}{, Pisa.}
\talkitem{Apr. 2024}{Higgs Centre}{Scattering and bound observables for spinning particles in Kerr}{, Edinburgh.}
\talkitem{Apr. 2024}{NTNU}{Gravitational bound waveforms from amplitudes}{, Taiwan, online.}
\talkitem{Mar. 2024}{ITMP, Moscow State University}{Gravitational bound waveforms from amplitudes}{, online.}
\talkitem{Feb. 2024}{Joint Belgian HEPTH Seminars}{Gravitational bound waveforms from amplitudes}{, Brussels.}
\talkitem{Nov. 2023}{Sapienza University of Rome}{Scattering and bound waveforms for Kerr black holes}{}
\talkitem{Nov. 2023}{University of Rome Tor Vergata}{Scattering and bound waveforms for Kerr black holes}{}
\talkitem{Nov. 2023}{STAG Centre}{Spinning waveforms from classical amplitudes}{, Southampton.}
\talkitem{May 2023}{IPhT}{Boundary-to-bound dictionary for generic Kerr orbits}{, Paris.}
\talkitem{Mar. 2023}{University of Geneva}{Shedding light on bound states with amplitudes}{}
\talkitem{Dec. 2022}{Higgs Centre}{Bethe--Salpeter equation for classical bound states}{, Edinburgh.}
\talkitem{Nov. 2022}{University of Nottingham}{Celestial holography on non-trivial backgrounds}{}
\talkitem{Apr. 2022}{University of Padova}{Amplitudes for the classical two-body problem in general relativity}{}
\talkitem{Mar. 2022}{University College Dublin}{Classical amplitudes for the two-body problem in general relativity}{}
\talkitem{Sept. 2021}{Caltech}{Gravitational event shapes and coherent states for gravitational waves}{}
\talkitem{June 2021}{Higgs Centre}{Coherent states from amplitudes and classical factorisation}{, Edinburgh.}
\talkitem{Feb. 2021}{NBI}{Waveforms from amplitudes and coherent states}{, Copenhagen.}
\talkitem{Jan. 2020}{Trinity College Dublin}{OPE at null infinity}{, Dublin.}
\talkitem{May 2017}{University of G\"ottingen}{The structure of intertwiners in the infinite-spin representation}{}
}

\section{Teaching activities}
Teaching and tutoring at high-school, Master's and PhD levels, primarily in English.\par\smallskip
\cvitem[0.19em]{2025}{\textbf{Invited Lecturer}, PhD school \emph{Analytic Computing in High-Energy and Gravitational Theoretical Physics}, Atrani. Delivered a series of PhD-level lectures entitled \emph{On-shell observables in the gravitational two-body problem} to about 40 graduate students and early-career researchers. \linktag{https://indico.cern.ch/event/1514553/}{school}}
\cvitem[0.19em]{2022}{\textbf{Tutor}, BUSSTEPP, Imperial College London. Two-week PhD-level summer school; tutorial and problem-solving sessions for about 80 international PhD students.}
\cvitem[0.19em]{2022}{\textbf{Tutor}, Conformal Bootstrap course, Higgs Centre School of Theoretical Physics. One-week PhD-level school course; tutorial and problem-solving sessions for about 50 international students.}
\cvitem[0.19em]{2022}{\textbf{Full-time teacher}, Differential Calculus, International School of Business / Niels Brock programme, Dublin. Three-month course for two high-school classes of about 40 students each; responsible for course delivery, regular assessment, final examination and grading.}
\cvitem[0.19em]{2019--2020}{\textbf{Teaching Assistant}, Trinity College Dublin. Master's-level courses in Quantum Field Theory II, Combinatorics and Algebraic Geometry (5 ECTS each); each course lasted about four months, with classes of about 30--40 students.}

\section{Academic service}
\cvitem[0.18em]{2022--}{Journal referee for \emph{Journal of High Energy Physics}, \emph{Physical Review D}, \emph{Physical Review Letters}, \emph{Journal of Cosmology and Astroparticle Physics} and \emph{SciPost Physics}.}
\cvitem[0.18em]{2025}{Member of a PhD selection committee at Queen Mary University of London.}
\cvitem[0.18em]{2024}{Panel member for MPhys thesis presentations in Astronomy and Cosmology at the University of Edinburgh; assessed about 50 candidates.}

\section{Scientific memberships}
\cvitem[0.18em]{2025--}{Core member of the Laser Interferometer Space Antenna (LISA) Consortium; initiating projects on EMRI waveform construction and validation, and on applications of particle-physics methods to gravitational-wave data analysis.}
\cvitem[0.18em]{2026--}{Core member of the Einstein Telescope (ET) Collaboration.}

\section{Research stays}
\cvitem[0.14em]{2024}{Humboldt University, Berlin (September--October); invited by Jan Plefka.}
\cvitem[0.14em]{2021}{Institut de Physique Th\'eorique (IPhT), Paris (November--December); invited by David A. Kosower.}
\cvitem[0.14em]{2020}{Higgs Centre for Theoretical Physics, Edinburgh (February); invited by Donal O'Connell.}
\cvitem[0.14em]{2019}{CERN, Geneva (August); invited by Claude Duhr.}

\section{Press and outreach}
\cvitem[0.18em]{2026}{Featured in the Nordita article \emph{Bringing together communities to improve precision in gravitational-wave modeling}, about the scientific aims of the programme and the experience of students and postdoctoral researchers who took part. \linktag{https://nordita.org/news-archive/news-2026/bringing-together-communities-to-improve-precision-in-gravitational-wave-modeling/}{article}}
\cvitem[0.18em]{2022}{Interviewed by Charlie Wood for the \emph{Quanta Magazine} article \emph{Massive Black Holes Shown to Act Like Quantum Particles}, about my work on coherent graviton states and the quantum-to-classical transition in black-hole scattering. \linktag{https://www.quantamagazine.org/massive-black-holes-shown-to-act-like-quantum-particles-20220329/}{article}}
\cvitem[0.18em]{2022}{Contributed personal video footage to the SAGEX film \emph{Doing a PhD in Physics}, presenting the daily life, challenges and rewards of doctoral research. \linktag{https://www.youtube.com/watch?v=ZbV0YnzeTzQ}{film}}
\cvitem[0.18em]{2021}{Contributed to the Feynman-diagram section of the SAGEX online exhibition \emph{At the Frontiers of Physics}, using short videos and interactive material to introduce particle physics to a broad audience. \linktag{https://exhibition.sagex.org/\#/exhibition-hub/}{exhibition}}
\cvitem[0.18em]{}{Public-engagement activities with \emph{Pint of Science} in Padova, Edinburgh and London.}

\section{Research software and external collaboration}
\cventry{2021}{Mathematics in Industry Study Group}{ESGI 165 / Defence Science and Technology Laboratory (DSTL)}{}{}{Contributed to the report \emph{Using the Electromagnetic Spectrum}. The project focused on estimating the location and power of an unknown transmitter from signal-strength measurements, using mathematical modelling, statistical inference and numerical methods. \linktag{https://www.cambridge.org/engage/api-gateway/miir/assets/orp/resource/item/61361b7c656369e082237850/original/using-the-electromagnetic-spectrum.pdf}{report}}
\cventry{2021}{Research Software Intern}{Wolfram Research}{Champaign, Illinois}{}{Developed code for the symbolic integration of products of Meijer $G$ functions with Devendra Kapadia. The implementation outperformed the standard \texttt{Integrate} function in a number of cases and was incorporated into Mathematica 14.0.}
\cvitem{Skills}{Mathematica, C++, \LaTeX, Git and Linux; experience with MATLAB and COMSOL Multiphysics.}
\cvitem{Languages}{Italian: mother tongue; English: fluent professional working language.}

\section{References}
\referenceitem{Prof. Ruth Britto}{Professor in Theoretical Physics, Trinity College Dublin}{britto@maths.tcd.ie}
\referenceitem{Prof. Donal O'Connell}{Professor of Theoretical Physics, University of Edinburgh}{donal@ed.ac.uk}
\referenceitem{Prof. Adam Pound}{Professor of Theoretical Physics, University of Southampton}{A.Pound@soton.ac.uk}
\par\noindent Additional references available upon request: Prof. Tim Adamo, Prof. Leor Barack, Prof. Vittorio Del Duca, Prof. David A. Kosower, Prof. Ira Z. Rothstein, Prof. Niels Warburton and Prof. Mao Zeng.

\end{document}
